%% Hand-authored circuitikz figure -- NOT generated by maestro2tex, placed by it via
%% the `order` list in the pixel block's maestro2tex config. Topology transcribed
%% instance-by-instance from the system bench netlist
%%   simruns/tb_anatop_pixel_iacc2/.../Ocean.0/psf/ps_iacc/netlist/input.scs
%% (cell castalia/tb_anatop_pixel, view PotFbTB). Supplies, En, the busset6 trim word
%% and the shared local bias generator are omitted -- they carry no measurement.
%% NOTE: the stb/noise tests (PotCtrlTB, PotTiaTB) disable VSRC_PATTERN and I_far and
%% drive V_InP from a single dc source V0 with the same dc value v_ref.
\begin{figure}[htbp]
  \centering
  \resizebox{\linewidth}{!}{%
\ctikzset{bipoles/length=1.0cm, resistors/scale=0.85, capacitors/scale=0.85, sources/scale=0.85}
\begin{circuitikz}[
    every node/.append style={font=\small},
    nlab/.style={font=\footnotesize, inner sep=1.5pt},
    opt/.style={font=\footnotesize, text=black!55, align=center, inner sep=2pt},
    railnode/.style={circle, fill=black, inner sep=1.05pt},
    prb/.style={draw, circle, minimum size=9pt, inner sep=0pt},
  ]
  %% ---------------------------------------------------------------- amplifiers
  %% circuitikz puts `-' on top and `+' on the bottom of an `op amp'; every branch
  %% that meets a pin is placed at that pin's height with (x,0 |- pin), so no run
  %% is ever diagonal.
  \node[op amp] (A) at (0,0) {};
  \node at (A.center) {\small A1};
  \node[op amp] (B) at (13.4,0) {};
  \node at (B.center) {\small A2};

  %% ------------------------------------------- potential program -> A1 non-inv
  \coordinate (VIP) at (-2.9,0 |- A.+);
  \draw (A.+) -- (VIP);
  \node[nlab, anchor=south] at (-1.95,0.06 |- A.+) {\texttt{V\_InP}};
  \draw (VIP) to[V, l_=$v_{\mathrm{step}}$] (-2.9,-1.75);
  \draw (-2.9,-1.75) to[V, l_=$v_{\mathrm{ref}}$] (-2.9,-3.4) node[ground]{};
  \node[nlab, anchor=west] at (-2.78,-1.75) {\texttt{net5}};
  \node[opt, anchor=north] at (-2.9,-3.9) {\texttt{VSRC\_PATTERN}\\ over \texttt{VSRC\_VCM}};

  %% ------------------------------ reference-electrode feedback through IPRB1
  \coordinate (VIN) at (-1.6,0 |- A.-);
  \draw (A.-) -- (VIN);
  \node[nlab, anchor=north] at (-1.05,-0.06 |- A.-) {\texttt{V\_InN}};
  \draw (VIN) -- (-1.6,2.6);
  \node[prb] (P1) at (0.8,2.6) {};
  \draw (-1.6,2.6) -- (P1.west);
  \node[nlab, above=2pt] at (P1.north) {\texttt{IPRB1}};
  \draw (P1.east) -- (4.8,2.6) -- (4.8,0);

  %% ------------------------------- A1 class-AB compensation, C1/C2 back to CE
  \draw (A.up) -- (A.up |- 0,1.5);
  \draw (A.up |- 0,1.5) to[C, l=$C_c$] (2.0,1.5) -- (2.0,0);
  \node[nlab, anchor=east, xshift=-2pt] at (A.up |- 0,1.08) {\texttt{C1}};
  \draw (A.down) -- (A.down |- 0,-1.5);
  \draw (A.down |- 0,-1.5) to[C, l_=$C_c$] (2.0,-1.5) -- (2.0,0);
  \node[nlab, anchor=east, xshift=-2pt] at (A.down |- 0,-1.08) {\texttt{C2}};

  %% --------------------------------------------------- three-electrode cell
  \draw[dashed, black!45] (2.25,-1.5) rectangle (9.75,1.5);
  \node[opt, anchor=north] at (6.0,-1.62) {\texttt{randles\_cell}};
  \draw (A.out) -- (2.0,0);
  \node[railnode] at (2.0,0){};
  \node[nlab, anchor=south east, xshift=-1pt, yshift=2pt] at (2.0,0) {\texttt{CE}};
  \draw (2.0,0) to[R, l=$R_{\mathrm{ce}}$] (4.8,0);
  \node[railnode] at (4.8,0){};
  \node[nlab, anchor=south east, xshift=-1pt, yshift=2pt] at (4.8,0) {\texttt{RE}};
  \draw (4.8,0) to[R, l=$R_u$] (7.2,0);
  \node[railnode] at (7.2,0){};
  \node[nlab, anchor=south east, inner sep=2.5pt] at (7.12,0.06) {$\mathrm{WE_s}$};
  \draw (7.2,0) -- (7.2,0.85) to[R, l=$R_{\mathrm{ct}}$] (9.5,0.85) -- (9.5,0);
  \draw (7.2,0) -- (7.2,-0.85) to[C, l_=$C_{\mathrm{dl}}$] (9.5,-0.85) -- (9.5,0);
  \draw (9.5,0) -- (10.0,0);

  %% --------------------------------------------- WE node, I_far, A2 inverting
  \node[railnode] at (10.0,0){};
  \node[nlab, anchor=north west, xshift=2pt, yshift=-2pt] at (10.0,0) {\texttt{WE}};
  \draw (10.0,0) to[I] (10.0,-3.4) node[ground]{};
  \node[nlab, anchor=east, xshift=-2pt] at (10.0,-2.5) {$i_{\mathrm{we}}$};
  \node[opt, anchor=north] at (10.0,-3.9) {\texttt{I\_far}};
  \draw (10.0,0) -- (10.0,2.4);
  \draw (B.-) -- (10.0,0 |- B.-);
  \node[railnode] at (10.0,0 |- B.-){};

  %% ------------------------------------------------------ V_CM at A2 non-inv
  \draw (B.+) -- (11.7,0 |- B.+);
  \draw (11.7,0 |- B.+) to[V, l_=$V_{\mathrm{CM}}$] (11.7,-3.4) node[ground]{};

  %% ---------------------------- A2 class-AB compensation, C1/C2 back to V_OUT
  \draw (B.up) -- (B.up |- 0,1.5);
  \draw (B.up |- 0,1.5) to[C, l=$C_c$] (15.0,1.5);
  \node[nlab, anchor=east, xshift=-2pt] at (B.up |- 0,1.08) {\texttt{C1}};
  \draw (B.down) -- (B.down |- 0,-1.5);
  \draw (B.down |- 0,-1.5) to[C, l_=$C_c$] (15.0,-1.5) -- (15.0,0);
  \node[nlab, anchor=east, xshift=-2pt] at (B.down |- 0,-1.08) {\texttt{C2}};

  %% ----------------------------------- transimpedance feedback through IPRB_F
  \node[prb] (PF) at (11.0,2.4) {};
  \draw (10.0,2.4) -- (PF.west);
  \node[nlab, above=2pt] at (PF.north) {\texttt{IPRB\_F}};
  \draw (PF.east) -- (12.2,2.4);
  \node[railnode] at (12.2,2.4){};
  \node[nlab, anchor=south, yshift=1pt] at (12.2,2.5) {\texttt{FB}};
  \draw (12.2,2.4) to[R, l=$R_f$] (15.0,2.4);
  \draw (15.0,2.4) -- (15.0,0);
  \node[railnode] at (15.0,1.5){};

  %% ------------------------------------------------------------ output node
  \draw (B.out) -- (15.0,0);
  \node[railnode] at (15.0,0){};
  \node[nlab, anchor=north west, xshift=2pt, yshift=-2pt] at (15.0,0) {\texttt{V\_OUT}};
  \draw (15.0,0) -- (16.2,0);
  \draw (16.2,0) to[C, l=$C_{\mathrm{out}}$] (16.2,-3.4) node[ground]{};
\end{circuitikz}
  }
  \caption[{System testbench for the potentiostat pixel.}]{The complete rev-2 potentiostat pixel as simulated, \texttt{tb\_anatop\_pixel} view \texttt{PotFbTB}: control amplifier A1 drives the counter electrode and closes its loop on the reference electrode, while transimpedance amplifier A2 holds the working electrode at $V_{\mathrm{CM}}$ and returns the cell current through $R_f$ to \texttt{V\_OUT}. Both amplifiers are the same \texttt{anatop\_tia\_amp\_ab} class-AB cell, each with its own external compensation capacitors $C_c = \SI{500}{\femto\farad}$ from pins \texttt{C1}/\texttt{C2} onto its own output; supplies, \texttt{En} and the shared local bias generator are omitted. \texttt{IPRB1} and \texttt{IPRB\_F} are \SI{0}{\volt} current probes inserted only so that \texttt{stb} can break the control loop and the transimpedance loop respectively -- they add no impedance and no offset. \texttt{I\_far} injects the accuracy-truth current with positive \texttt{iwe} flowing from the working electrode to ground, so $v_{\mathrm{out}} = V_{\mathrm{CM}} + \texttt{iwe}\cdot R_f$ and the measured current is the probe current in \texttt{IPRB\_F}; the opposite source orientation inverts the sign of every accuracy number. Cell elements are $R_{\mathrm{ce}} = R_u = \SI{1}{\kilo\ohm}$ fixed with $R_{\mathrm{ct}}$ and $C_{\mathrm{dl}}$ swept over the electrode grid, $R_f$ is the programmable feedback resistor (\SIrange{35}{560}{\kilo\ohm}), and $C_{\mathrm{out}} = \SI{5}{\pico\farad}$.}
  \label{fig:tb-pixel}
\end{figure}
